% =====================================================================
%  02 · Background — The Approaches
% =====================================================================
\section{Background \& Approaches}

\begin{frame}{Variational quantum circuits in one slide}{Expressivity vs.\ trainability}
    \begin{columns}[T]
        \begin{column}{0.5\textwidth}
            \begin{block}{What makes them expressive}
                \textbf{Data re-uploading}~\cite{perezsalinas2020reuploading}: the
                input is re-injected at every layer. Even a single qubit can become
                a universal classifier. A VQC is effectively a \textbf{Fourier
                    series} in its input~\cite{schuld2021encoding}.
            \end{block}
        \end{column}
        \begin{column}{0.5\textwidth}
            \begin{alertblock}{What makes them hard}
                \textbf{Barren plateaus}~\cite{mcclean2018barren}: in deep, wide,
                random circuits gradients vanish \emph{exponentially} in the qubit
                count --- and training stalls.
            \end{alertblock}
        \end{column}
    \end{columns}
    \vspace{2pt}
    \begin{callout}
        Width, depth and trainability trade off against each other --- a tension our
        architecture sweeps probe directly.
    \end{callout}
\end{frame}

\begin{frame}{Four quantum paradigms --- one fair test each}{Roadmap of the talk}
    \begin{columns}[T]
        \begin{column}{0.5\textwidth}
            \begin{block}{1 · Trainable circuits (VQC)}
                A small variational circuit as a classifier \textbf{head}.
                \\[2pt]\footnotesize Exp 1 (binary) · Exp 2 (4-class, frozen features)
            \end{block}
            \begin{block}{2 · Quantum kernels (QSVM)}
                The circuit as a fixed \textbf{feature
                    map}~\cite{havlicek2019,schuld2019feature} for a support-vector
                machine.
                \\[2pt]\footnotesize Exp 3 (few-shot)
            \end{block}
        \end{column}
        \begin{column}{0.5\textwidth}
            \begin{block}{3 · Quanvolution}
                A small circuit \textbf{slid across} the signal like a convolutional
                kernel~\cite{henderson2020quanvolutional}.
                \\[2pt]\footnotesize Exp 4 (noise robustness)
            \end{block}
            \begin{block}{4 · True QCNN}
                The whole spectrum embedded into \textbf{one quantum
                    state}~\cite{cong2019qcnn}, with coherent pooling.
                \\[2pt]\footnotesize Exp 5 (fully quantum)
            \end{block}
        \end{column}
    \end{columns}
    \vspace{2pt}
    \centering{\small\color{zhawgrey}Every paradigm is benchmarked against a
        parameter-matched classical control.}
\end{frame}
